\documentclass[11pt,a4paper]{article}

\usepackage{amsmath,amssymb,amsthm}
\usepackage{geometry}
\usepackage{xcolor}

\geometry{margin=2.5cm}

\theoremstyle{plain}
\newtheorem{theorem}{Theorem}
\newtheorem{corollary}[theorem]{Corollary}
\newtheorem{proposition}[theorem]{Proposition}
\theoremstyle{definition}
\newtheorem{definition}{Definition}
\newtheorem{example}{Example}

\newcommand{\Freq}{\mathrm{Freq}}
\newcommand{\freq}{\mathrm{freq}}
\newcommand{\GG}{\mathcal{G}}
\newcommand{\PP}{\mathcal{P}}

\newcommand\subsetsim{\mathrel{\substack{
  \textstyle\subset\\[-0.2ex]\textstyle\sim}}}

\title{Relation Between Supergraph-Based Frequency ($\Freq_G$)\\and Exact-Match Frequency ($\freq_{exact}$)}
\author{}
\date{}

\begin{document}
\maketitle

\section{Definitions}

Let $G = (V, E)$ be a graph. Let $\GG(G)$ denote the set of all connected subgraphs of~$G$.

\begin{definition}[Exact-match frequency]
For a pattern $P$, the \textbf{exact-match frequency} counts the number of distinct connected subgraphs of $G$ that are isomorphic to $P$:
\[
\freq_{exact}(P) \;=\; \big|\{H \in \GG(G) : H \cong P\}\big|
\]
\end{definition}

\begin{definition}[Supergraph-based frequency (proposed in the paper)]
The \textbf{supergraph-based frequency} counts the number of connected subgraphs of $G$ that \emph{contain} an isomorphic copy of $P$:
\[
\Freq_G(P) \;=\; |UP_G(P)| \;=\; \big|\{H \in \GG(G) : P \subsetsim H\}\big|
\]
where $P \subsetsim H$ means $\exists\, H' \subseteq H$ such that $H' \cong P$ (subgraph isomorphism).
\end{definition}

\section{Exact Relation}

\begin{theorem}[Summation over the subgraph lattice]\label{thm:relation}
For any pattern $P$ and any graph $G$:
\[
\boxed{\;\Freq_G(P) \;=\; \sum_{\substack{Q \,\in\, \GG(G)/{\cong} \\[2pt] P \,\subsetsim\, Q}} \freq_{exact}(Q)\;}
\]
where $\GG(G)/{\cong}$ is the set of isomorphism classes of connected subgraphs of~$G$. Equivalently:
\[
\Freq_G(P) \;=\; \freq_{exact}(P) \;+\; \sum_{\substack{Q \,\in\, \GG(G)/{\cong} \\[2pt] P \,\subsetneq_i\, Q}} \freq_{exact}(Q)
\]
where $P \subsetneq_i Q$ means $P$ is isomorphic to a \emph{proper} subgraph of $Q$.
\end{theorem}

\begin{proof}
Every subgraph $H \in \GG(G)$ belongs to exactly one isomorphism class $[H] \in \GG(G)/{\cong}$. We can therefore partition $UP_G(P)$ by the isomorphism class of its elements:
\begin{align*}
\Freq_G(P) \;&=\; \big|\{H \in \GG(G) : P \subsetsim H\}\big| \\[4pt]
             &=\; \sum_{H \in \GG(G)} \mathbf{1}[P \subsetsim H] \\[4pt]
             &=\; \sum_{Q \,\in\, \GG(G)/{\cong}} \;\sum_{\substack{H \in \GG(G) \\[1pt] H \cong Q}} \mathbf{1}[P \subsetsim H] \\[4pt]
             &=\; \sum_{Q \,\in\, \GG(G)/{\cong}} \mathbf{1}[P \subsetsim Q] \cdot \big|\{H \in \GG(G) : H \cong Q\}\big| \\[4pt]
             &=\; \sum_{\substack{Q \,\in\, \GG(G)/{\cong} \\[2pt] P \,\subsetsim\, Q}} \freq_{exact}(Q).
\end{align*}
The key step (line 3 to 4) uses the fact that $P \subsetsim H$ depends only on the isomorphism class of $H$: if $H \cong Q$ and $P \subsetsim Q$, then $P \subsetsim H$.

Separating the term $Q = P$ (where $P \subsetsim P$ trivially via the identity) gives the second form.
\end{proof}

\section{Interpretation}

\subsection{$\Freq_G$ as a cumulative sum over the subgraph lattice}

The set of isomorphism classes $\GG(G)/{\cong}$ forms a partially ordered set (poset) under $\subsetsim$. The function $\freq_{exact}$ assigns a count to each element of this poset, and $\Freq_G$ is its \textbf{cumulative sum upward} (the zeta transform):
\[
\Freq_G(P) = \big(\zeta \cdot \freq_{exact}\big)(P) \;=\; \sum_{Q \,\geq\, P} \freq_{exact}(Q)
\]
where $Q \geq P$ means $P \subsetsim Q$ in the poset ordering. This is a standard construction in combinatorics on partially ordered sets.

\subsection{Recovering $\freq_{exact}$ from $\Freq_G$ (M\"obius inversion)}

By M\"obius inversion on the poset $(\GG(G)/{\cong},\; \subsetsim)$, the exact frequency can be recovered:
\[
\freq_{exact}(P) \;=\; \sum_{\substack{Q \,\in\, \GG(G)/{\cong} \\[2pt] P \,\subsetsim\, Q}} \mu(P, Q) \cdot \Freq_G(Q)
\]
where $\mu$ is the M\"obius function of the subgraph poset, defined recursively by $\mu(P, P) = 1$ and
\[
\mu(P, Q) = -\sum_{\substack{R \\[1pt] P \,\subsetsim\, R \,\subsetneq_i\, Q}} \mu(P, R) \quad \text{for } P \subsetneq_i Q.
\]
Computing $\mu$ is intractable in general (the subgraph lattice is exponentially large), but this shows that $\Freq_G$ and $\freq_{exact}$ are \emph{in bijection}: knowing one determines the other completely.

\subsection{Why $\Freq_G$ is the right measure for anomaly detection}

The relation $\Freq_G(P) = \freq_{exact}(P) + \sum_{Q \supsetneq P} \freq_{exact}(Q)$ means:

\begin{itemize}
\item $\freq_{exact}(P)$ counts the occurrences of $P$ \emph{alone} --- how many times $P$ appears as an exact copy.
\item $\Freq_G(P)$ counts the occurrences of $P$ \emph{in any context} --- how many subgraphs of $G$ contain $P$ as part of their structure, regardless of what else they contain.
\end{itemize}

For anomaly detection, $\Freq_G$ is more appropriate because:
\begin{enumerate}
\item A pattern $P$ that is rare as an exact match but common inside larger structures is \emph{not} anomalous --- it is a building block of normal behavior. Only $\Freq_G(P)$ captures this: if many larger subgraphs contain $P$, then $\Freq_G(P)$ is high, correctly indicating normality.

\item Conversely, if $\Freq_G(P)$ is low, then $P$ is rare both as an exact match \emph{and} as a substructure of any larger pattern. This is a stronger (and more useful) statement of anomalousness.
\end{enumerate}

\section{Bounds}

\begin{proposition}[Bounds on $\Freq_G$ in terms of $\freq_{exact}$]
\label{prop:bounds}
For any pattern $P$:
\[
\freq_{exact}(P) \;\leq\; \Freq_G(P) \;\leq\; \sum_{s \,=\, |V(P)|}^{n} N_s(G)
\]
where $N_s(G) = |\{H \in \GG(G) : |V(H)| = s\}|$ is the number of connected subgraphs of $G$ with exactly $s$ nodes.

The lower bound is tight when $P$ is never contained in any strictly larger connected subgraph of $G$ (e.g., $P = G$ itself). The upper bound is tight when $P$ is a single node or edge (contained in every connected subgraph of sufficient size).
\end{proposition}

\begin{proposition}[Anti-monotonicity from the summation identity]
If $P_1 \subseteq_i P_2$, then $\{Q : P_2 \subsetsim Q\} \subseteq \{Q : P_1 \subsetsim Q\}$ by transitivity of $\subsetsim$. Therefore:
\[
\Freq_G(P_2) = \sum_{Q \,\supseteq\, P_2} \freq_{exact}(Q) \;\leq\; \sum_{Q \,\supseteq\, P_1} \freq_{exact}(Q) = \Freq_G(P_1)
\]
since the sum for $P_2$ runs over a subset of the terms in the sum for $P_1$, and all terms $\freq_{exact}(Q) \geq 0$.
\end{proposition}

\section{Examples}

\begin{example}[Single edge]
Let $P = K_2$ (a single edge). Then $P \subsetsim H$ for every connected subgraph $H$ with $|V(H)| \geq 2$. So:
\[
\Freq_G(K_2) = \sum_{\substack{Q \\[1pt] |V(Q)| \geq 2}} \freq_{exact}(Q) = |\GG(G)| - n
\]
since the only connected subgraphs NOT containing an edge are the $n$ isolated nodes. Meanwhile $\freq_{exact}(K_2) = |E|$.
\end{example}

\begin{example}[Triangle]
Let $P = K_3$ (a triangle). Then:
\[
\Freq_G(K_3) = \freq_{exact}(K_3) + \sum_{\substack{Q \supsetneq K_3}} \freq_{exact}(Q)
\]
The first term counts the triangles in $G$. The second counts every larger connected subgraph of $G$ that contains at least one triangle. In dense graphs, $\Freq_G(K_3) \gg \freq_{exact}(K_3)$.
\end{example}

\begin{example}[Anomalous clique]
Let $P = K_7$ (a 7-clique injected as an anomaly into a sparse graph like Cora with average degree 4.1). Then:
\begin{itemize}
    \item $\freq_{exact}(K_7)$: the number of 7-cliques in $G$. In Cora with 10 injected cliques, $\freq_{exact}(K_7) = 10$.
    \item $\Freq_G(K_7) = 10 + \sum_{Q \supsetneq K_7} \freq_{exact}(Q)$. Since Cora is sparse, there are very few connected subgraphs larger than 7 nodes that still contain a complete 7-clique. So $\Freq_G(K_7) \approx \freq_{exact}(K_7) = 10$. The two measures nearly coincide for large, rare patterns.
\end{itemize}
This illustrates that $\Freq_G(P) \approx \freq_{exact}(P)$ when $P$ is large and rare --- precisely the regime where anomaly detection operates.
\end{example}

\section{Summary}

The exact, deterministic relation between the two frequency measures is:
\[
\boxed{\;\Freq_G(P) = \sum_{\substack{Q \,\in\, \GG(G)/{\cong} \\[2pt] P \,\subsetsim\, Q}} \freq_{exact}(Q) = \freq_{exact}(P) + \sum_{\substack{Q \,\supsetneq\, P}} \freq_{exact}(Q)\;}
\]

$\Freq_G$ is the zeta transform of $\freq_{exact}$ over the subgraph lattice, and $\freq_{exact}$ is recoverable via M\"obius inversion. For anomaly detection, where the patterns of interest are large and rare, the surplus $\sum_{Q \supsetneq P} \freq_{exact}(Q)$ is small, making $\Freq_G(P) \approx \freq_{exact}(P)$. For small, common patterns, the surplus is large, but this is desirable: it correctly prevents flagging common building blocks as anomalous.

\end{document}
